% Generated draft: requires mathematical and source review.
\documentclass[11pt]{article}
\usepackage[margin=25mm]{geometry}
\usepackage{fontspec}
\setmainfont{Latin Modern Roman}
\usepackage{amsmath,amssymb,mathtools,unicode-math}
\setmathfont{Latin Modern Math}
\usepackage{enumitem,needspace,xurl,hyperref}
\setlength{\parindent}{0pt}
\setlength{\parskip}{5pt}
\setlength{\emergencystretch}{3em}
\newcommand{\N}{\mathbb{N}}
\newcommand{\Z}{\mathbb{Z}}
\newcommand{\Q}{\mathbb{Q}}
\newcommand{\R}{\mathbb{R}}
\newcommand{\C}{\mathbb{C}}
\newcommand{\F}{\mathbb{F}}
\newcommand{\A}{\mathbb{A}}
\newcommand{\PP}{\mathbb{P}}
\newcommand{\E}{\mathbb{E}}
\newcommand{\eps}{\varepsilon}
\newcommand{\dd}{\,\mathrm{d}}
\begin{document}
\section*{UnsolvedMath ID: 20000939}
\textbf{Source problem code:} catalogue entry
\subsection{Definitions and mathematical statement}
The source's semi-algebraic Ramsey question, attributed to Bukh--Matou\v{s}ek, asks whether an omitted preceding bound can be improved to a constant number of iterated logarithms. The extraction does not define $f$, $r_f$ or the antecedent ``this.'' A precise natural interpretation fixes a semi-algebraic $k$-ary predicate $f$ on ordered distinct points of $\R^d$, expressed by a Boolean combination of polynomial sign conditions, and asks for a homogeneous subset bound whose number of logarithms is independent of $k$ and $d$ under specified control of the predicate.

Let $R_f^{(k)}(s,n)$ be the least $N$ such that every ordered $N$-point set contains either $s$ points whose increasing $k$-tuples all satisfy $f$, or $n$ whose increasing $k$-tuples all fail $f$. Define $\operatorname{tw}_1(x)=x$ and $\operatorname{tw}_{h+1}(x)=2^{\operatorname{tw}_h(x)}$. This explicitly supplies conventions absent from the fragment.

The attached aid reports that an Adibelli--Tomon February 2026 preprint proves, for total polynomial degree bounded by an integer $D\ge1$,
\[
R_f^{(k)}(n,n)\le\operatorname{tw}_{3D^3}(C n),
\]
with tower height independent of uniformity and dimension; other dependence of $C$ and the precise complexity hypotheses require the preprint statement. Inverting the tower gives at most $3D^3-1$ iterated logarithms for a guaranteed homogeneous size. This answers fixed-predicate or bounded-degree readings as reported, while the stronger dimension-only question uniform over degrees is expressly not settled.

For its refinement, divide variables into $k$ point-coordinate blocks and define the monomial block-support width $w$ as the largest number of distinct point blocks occurring in a monomial of any defining polynomial. The aid states that, for $1\le w\le k$, its quasi-algebraic theorem yields
\[
R_f^{(k)}(n,n)\le\operatorname{tw}_{K_w}(C n),\qquad
K_w=\frac{w(w^2+6w+1)}2,
\]
independent of the total powers within participating blocks. It calls the resulting relation $w$-dependent, but the exact quasi-algebraic definition is omitted. For $w=1$ it reports an injective coordinate lift and a stronger Jin--Tomon semi-linear one-exponential bound, without supplying that bound's complete parameter dependence. For $w=0$ all predicates are constant and every vertex set is homogeneous.

The original context attributes to Bukh--Matou\v{s}ek a sharp two-logarithm result in dimension one. It attributes to Conlon--Fox--Pach--Sudakov--Suk the written bound $r_f(N)\ge c\log^{(k-2)}N$, where $\log^{(j)}$ means $j$ iterated logarithms, and sharpness in sufficiently high dimension $\R^{c(k)}$. Because the definition and normalization of $r_f$ are missing, this literal statement's possible power losses cannot be checked here. The original also asks for the unbalanced variant $R_f^{(k)}(s,n)$; the aid's displayed improvements concern diagonal thresholds only.
\subsection{Short English statement}
Clarify how many iterated logarithms are needed for semi-algebraic homogeneous sets. The supplied aid reports a preprint answer under bounded degree and a block-support refinement, while leaving the degree-uniform dimension-only question and unbalanced variant unresolved.
\subsection{Sources}
\begin{itemize}
\item Local UnsolvedMath registry, permanent ID $20000939$, source code catalogue entry; canonical \texttt{aim-combinatorics-notes.json}, \texttt{notes[63]}, tagged question.
\item AIM ``Graph Ramsey theory,'' ``Semi-algebraic Ramsey numbers,'' source item 17: \url{https://www.overleaf.com/read/mnvcscjjysvg}. The source names Bukh--Matou\v{s}ek and Conlon--Fox--Pach--Sudakov--Suk without bibliographic identifiers; the fragment's missing antecedent is not supplied.
\item Attached research attempt 1: partial result, run and registry status partially solved, rigor proved, novelty confidence low, classification PARTIAL-PROGRESS. It names an Adibelli--Tomon February 2026 preprint and Jin--Tomon, without exact citations. These machine-generated, non-peer-reviewed claims, including the preprint interpretation, require independent verification before citation.
\end{itemize}
\end{document}
